
\medskip
\phantomsection\label{par:sga3-VIA-3.2.1}%
\noindent\textbf{3.2.1.---}
Suppose for a moment that every point of \(G\) has an open saturated
neighborhood \(W\) such that the groupoid induced by \(G_*\) on \(W\)
has a quasi-section (Expos\'e~V, \S~6). Theorem~3.2 then follows from
\SGARefLink{lem:sga3-V-6.1}{Expos\'e~V, 6.1}. We shall show that, moreover, \emph{every finite
subset of \(F\backslash G\) is then contained in an affine open subset}.\footnote{Translator's note: the Gille--Polo re-edition
(2011), Exposé~VI\textsubscript{A}, adds that when \(A=k\) is a field every quasi-compact open subset
of \(F\backslash G\) is even quasi-projective (a result of W.-L.~Chow for smooth groups), referring to
M.~Raynaud, \emph{Faisceaux amples sur les schémas en groupes et les espaces homogènes}, Lecture
Notes in Math.\ \textbf{119}, Springer, 1970, VI, 2.6; and that over \(A=\mathbb C[\varepsilon]/(\varepsilon^2)\)
there are abelian \(A\)-schemes that are not projective (loc.\ cit., XII, 4.2).}

If \(U\) is a quasi-section of the groupoid induced by \(G_*\) on a
saturated open subset \(W\) of \(G\), then \(U\otimes_AA'\) is indeed a
quasi-section of the \((\mathrm{Sch}/A')\)-groupoid induced by
\(G_*\otimes_AA'\) on \(W\otimes_AA'\). Moreover, if \(U_*\) is the
\((\mathrm{Sch}/A)\)-groupoid induced by \(G_*\) on \(U\), then
\(U_*\otimes_AA'\) is identified with the
\((\mathrm{Sch}/A')\)-groupoid induced by \(G_*\otimes_AA'\) on
\(U\otimes_AA'\). It therefore follows from the proofs of Expos\'e~V that the
construction of the quotient \(F\backslash G\) commutes with a base extension
\(A\to A'\) of the kind considered here.\footnote{Translator's note: the point is that the formation of
the direct images of the structure sheaves by the morphisms of the construction (the projection
onto the quotient, \(\lambda\) and \(\operatorname{pr}_2\)) commutes with the flat base change
\(A\to A'\); as \(F\) and \(G\) are of finite type over \(A\), these morphisms are quasi-compact and
quasi-separated, so this is flat base change for quasi-coherent direct images (EGA~III\textsubscript{1},
1.4.15; EGA~IV\textsubscript{1}, 1.7.21). See also 4.6 below. The Gille--Polo re-edition (2011),
Exposé~VI\textsubscript{A}, adds this explanation.}

Let then \(x_1,\ldots,x_n\) be points of \(F\backslash G\), which we may
suppose closed,\footnote{Translator's note: \(F\backslash G\) is here of finite type over \(A\), so every
point specializes to a closed point, and an open subset containing that closed point contains the
point; the Gille--Polo re-edition (2011), Exposé~VI\textsubscript{A}, spells this out.} and let \(g_1,\ldots,g_n\) be closed points of \(G\) mapping to
\(x_1,\ldots,x_n\); let \(U\) be an everywhere dense affine open subset of
\(F\backslash G\),\footnote{Translator's note: such an open subset exists because \(F\backslash G\) is
noetherian: in each of its finitely many irreducible components, choose a nonempty affine open
subset of \(F\backslash G\) that does not meet the other components; these are pairwise disjoint, and
their union is affine and dense. The Gille--Polo re-edition (2011), Exposé~VI\textsubscript{A}, explains
this.} \(V\) its inverse image in \(G\), and \(A'\) an \(A\)-algebra
sufficiently large for our purposes. Then the inverse image of
\(g_1,\ldots,g_n\) in \(G\otimes_AA'\) consists of points
\(g'_1,\ldots,g'_p\) which are strictly rational over \(A'\)
(\hyperref[prop:sga3-VIA-1.2]{1.2}); moreover, if \(A'\) is sufficiently large, the everywhere
dense open subset
\[
  \bigcap_i\,(V\otimes_AA')^{-1}\cdot g'_i
\]
contains a point \(x\) strictly rational over \(A'\).\footnote{Translator's note: this open subset is
not empty, hence contains a closed point; by 1.2, this point becomes strictly rational after a further
local, finite and free extension of \(A'\), and the open subset and the points \(g'_i\) are
compatible with such an extension (note to 0.4). The Gille--Polo re-edition (2011),
Exposé~VI\textsubscript{A}, details this step.} Hence
\(x\in(V\otimes_AA')^{-1}\cdot g'_i\), whence
\(g'_i\in(V\otimes_AA')\cdot x\). Since \((V\otimes_AA')\cdot x\) is the image
of \(V\otimes_AA'\) under an automorphism of the groupoid \(G_*\otimes_AA'\), the
image \((U\otimes_AA')\cdot x\) of \((V\otimes_AA')\cdot x\) in
\((F\backslash G)\otimes_AA'\) is the image of \(U\otimes_AA'\) under the
automorphism induced on \((F\backslash G)\otimes_AA'\). Therefore
\((U\otimes_AA')\cdot x\) is an affine open subset of
\((F\backslash G)\otimes_AA'\) containing the images of \(g'_1,\ldots,g'_p\).

To complete the proof, one shows as in \SGARefLink{par:sga3-V-5-b}{Expos\'e~V, \S~5~b)} that
\((U\otimes_AA')\cdot x\) contains an affine open subset \(W'\) which contains
the images of \(g'_1,\ldots,g'_p\) and is saturated for the equivalence relation
\begin{center}
  \phantomsection\label{diag:sga3-VIA-3.2.1-descent-equivalence}%
  \begin{tikzcd}[column sep=large]
    (F\backslash G)\otimes_A A'\otimes_A A'
      \arrow[r,shift left=1.2ex]
      \arrow[r,shift right=1.2ex]
    &
    (F\backslash G)\otimes_A A'
  \end{tikzcd}
\end{center}
induced by the two canonical injections of \(A'\) into \(A'\otimes_AA'\). The
image of \(W'\) in \(F\backslash G\) then contains \(x_1,\ldots,x_n\) and is an
affine open subset of \(F\backslash G\) (\SGARefLink{thm:sga3-V-4.1}{Expos\'e~V, 4.1}).
